% ============================================================
%  Capítulo: Metodología
% ============================================================

\chapter{Metodología}

\section{Métodos}
\subsection{Datos y generación de vistas}

El conjunto de datos utilizado corresponde a un subconjunto proveniente de
ShapeNet~\cite{chang2015shapenet}, compuesto por 1.700 objetos divididos en
dos categorías de simetría: 850 objetos con simetría axial (\texttt{axis\_sym})
y 850 con simetría planar (\texttt{plane\_sym}) (ver Figura~\ref{fig:tipos_simetria}).

Cada objeto se representa mediante una malla poligonal en formato \texttt{.obj},
centrada en el origen y normalizada a escala unitaria, junto con un archivo de
anotación \texttt{.txt} que contiene los parámetros de referencia de su simetría.
Para objetos axiales se especifica un vector de dirección $\hat{\mathbf{d}}$ y un
punto de origen $\mathbf{o}$ del eje, mientras que para los objetos planares se
especifican entre uno y tres planos, cada uno definido por su vector normal
$\hat{\mathbf{n}}$ y un punto de origen $\mathbf{o}$.

\begin{figure}[ht]
    \centering
    \begin{subfigure}{0.3\linewidth}
        \centering
        \includegraphics[width=\linewidth]{axis_sym_no_visual.png}
        \caption{Objeto con simetría axial, simetría no visualizable.}
        \label{fig:axis_no_visual}
    \end{subfigure}
    \hfill
    \begin{subfigure}{0.3\linewidth}
        \centering
        \includegraphics[width=\linewidth]{axys_sym_visual.png}
        \caption{Objeto con simetría axial, simetría visualizable.}
        \label{fig:axis_visual}
    \end{subfigure}

    \vspace{0.5cm}

    \begin{subfigure}{0.3\linewidth}
        \centering
        \includegraphics[width=\linewidth]{planar_sym_visual.png}
        \caption{Objeto con simetría planar, simetría no visualizable.}
        \label{fig:planar_no_visual}
    \end{subfigure}
    \hfill
    \begin{subfigure}{0.3\linewidth}
        \centering
        \includegraphics[width=\linewidth]{planar_sym_no_visual.png}
        \caption{Objeto con simetría planar, simetría visualizable.}
        \label{fig:planar_visual}
    \end{subfigure}

    \caption{Ejemplos de objetos con simetría axial y planar.}
    \label{fig:tipos_simetria}
\end{figure}

\newpage

Para cada uno de los objetos se genera un conjunto de 114 vistas mediante
renderizado multivista basado en muestreo de Fibonacci sobre la superficie de
una esfera (véase Figura~\ref{fig:fibonacci}). Este método distribuye los puntos
de vista de manera aproximadamente uniforme, evitando concentraciones angulares
y, por tanto, sesgos en la cobertura del objeto.

Cada punto de la esfera define la posición de una cámara virtual orientada hacia
el origen, con un campo de visión de $\text{FoV} = 60°$. El renderizado se
realiza mediante PyTorch3D empleando una cámara de perspectiva
(\texttt{FoVPerspectiveCamera}), y para cada vista se almacenan en
\texttt{metadata\_all.json}:

\begin{itemize}
    \item la matriz de rotación $R \in \mathbb{R}^{3\times 3}$ (convenio
          PyTorch3D de vector-fila: $p_{\text{cam}} = p_{\text{world}}\,R + T$),
    \item el vector de traslación $T \in \mathbb{R}^{3}$,
    \item la posición de la cámara en coordenadas mundiales $\mathbf{C} = -R\,T$,
    \item los ángulos de azimut y elevación correspondientes.
\end{itemize}

Las variables experimentales consideradas en esta etapa son:

\begin{itemize}
    \item \textbf{Número de vistas} ($N$): configuraciones de 1, 6, 14 y 26
          vistas para la evaluación base, con extensiones a 42, 62, 86 y 114
          vistas para análisis de saturación.
    \item \textbf{Resolución de imagen}: 224, 448 y 1136 píxeles, para evaluar
          el efecto sobre la precisión de la retroproyección 2D$\rightarrow$3D.
    \item \textbf{Iluminación}: tres condiciones de iluminación uniforme
          (\textit{flat}, valor de gris 0,7; \textit{darker}, 0,3;
          \textit{brighter}, 0,95), que permiten evaluar la robustez del modelo
          ante variaciones en el contraste sin introducir complejidad adicional
          por sombras direccionales.
\end{itemize}

Para garantizar reproducibilidad y ejecución incremental, cada combinación de
objeto, tamaño e iluminación se almacena en una estructura de directorios
independiente junto con \texttt{metadata\_all.json}.

\newpage
\subsection{Extracción de puntos geométricos mediante VLM}
\label{sec:extraccion_vlm}

Cada conjunto de vistas se procesa mediante \textit{Molmo2-8B}~\cite{clark2026molmo2},
que opera como una caja negra sin \textit{fine-tuning} adicional para esta tarea.
El modelo recibe las $N$ imágenes de un objeto junto con el \textit{prompt}
correspondiente y devuelve coordenadas 2D en un formato estructurado.

Las coordenadas de Molmo2 utilizan una escala normalizada de $[0, 1000]$ con
origen en la esquina superior izquierda de la imagen, independientemente de la
resolución efectiva de la imagen enviada. El formato de salida del modelo es:

\begin{equation}
    \texttt{<points\ coords="$f$\ $o$\ $X$\ $Y$;\ \ldots">\ \textit{descripción}\ </points>}
\end{equation}

donde los valores están separados por espacios y los puntos correspondientes a
imágenes distintas se separan por punto y coma. Cada punto queda especificado por
cuatro campos:

\begin{itemize}
    \item $f$: índice de la imagen (\textit{frame}) dentro del conjunto enviado,
          comenzando en 1.
    \item $o$: índice de objeto (\textit{object index}), único por objeto distinto
          señalado. Los índices son secuenciales a partir de 1, de modo que el
          último índice representa el conteo total de objetos señalados.
    \item $X, Y$: coordenadas normalizadas en $[0, 1000]$.
\end{itemize}

Los puntos se ordenan primero por índice de imagen y luego por coordenadas $X, Y$.
Para el caso multivista, la salida incluye un par de puntos por imagen, con
\texttt{obj\_id} $\in \{1, 2\}$ identificando cada uno. A modo de ejemplo, la
siguiente salida indica que en la imagen~1 se señaló el objeto en $(555, 169)$,
en la imagen~2 se señalaron dos objetos distintos en $(649, 154)$ y $(709, 162)$,
y en la imagen~5 se señalaron tres:

\begin{verbatim}
<points coords="1 1 555 169;2 1 649 154 2 709 162;5 1 758 175 2 808 183 3 852 187">
\end{verbatim}

Este formato compacto —en lugar de JSON u otros esquemas estructurados— reduce
significativamente el número de \textit{tokens} necesarios para representar las
predicciones~\cite{clark2026molmo2}.

\subsubsection{Estrategias de prompt (Flujo A — pointing directo)}

En la estrategia base, el modelo recibe directamente las $N$ vistas del objeto
junto con un prompt que especifica qué puntos estructurales señalar. Para obtener
los prompts con mejor rendimiento se diseñaron pruebas iterativas sobre conjuntos
de 100 objetos axiales y planares, cuyo detalle se explica en la
sección~\ref{Experimentos}.

\paragraph{Variantes de prompt para simetría axial:}
proyección directa del eje (puntos sobre la proyección del eje cerca de los
extremos superior e inferior), pares bilaterales (un punto y su espejo respecto
al eje), extremos de la silueta más ancha, pares de características estructurales
simétricas (asas, orificios, nervaduras), polos del eje de rotación, y puntos
sobre la línea central del eje.

\paragraph{Variantes de prompt para simetría planar:}
traza del plano sobre la superficie, pares bilaterales respecto al plano, traza
visible de la ``costura'' entre las dos mitades especulares, pares de elementos
estructurales correspondientes (patas, ruedas, brazos), puntos medios horizontales
de la silueta, y extremos de la traza visible del plano.

Cada variante incluye instrucciones explícitas de consistencia global a través de
todas las vistas, solicitando al modelo que infiera una hipótesis de simetría
conjunta antes de proyectar los puntos individualmente en cada imagen.

\subsubsection{Descripción del objeto como contexto (Flujo B — pointing con descripción)}

Una limitación del Flujo A es que el modelo recibe únicamente instrucciones
geométricas abstractas sin ningún conocimiento previo del tipo de objeto que está
analizando. Para mitigar esto, se propone un flujo aumentado en el que una
descripción semántica del objeto se incorpora al prompt antes de las instrucciones
de pointing:

\begin{enumerate}
    \item \textbf{Selección de vista de descripción}: se selecciona una vista
          representativa de las 114 disponibles (criterio detallado en la
          sección~\ref{sec:seleccion_vista}).
    \item \textbf{Generación de descripción}: esa vista única se envía a Molmo2
          con un prompt de descripción libre (\textit{``Describe briefly what
          object this is and its main geometric structure''}).
    \item \textbf{Prompt aumentado}: la descripción generada se prepende al prompt
          de pointing habitual, de modo que el modelo recibe las $N$ vistas junto
          con el contexto semántico del objeto.
\end{enumerate}

La hipótesis es que el contexto semántico permite al modelo orientar mejor sus
predicciones hacia las características estructuralmente relevantes del objeto
específico (e.g., saber que es una silla orienta los puntos hacia el respaldo y
patas, no hacia el borde del asiento).

\subsubsection{Descripción y pointing integrados (Flujo C)}

Como tercer flujo se propone mantener la solicitud de una descripción del objeto
mediante una vista de la 114, sin embargo el prompt a usar es distinto al del
Flujo A, y en cambio, se busca detectar todos los puntos señalados como
importantes:

\begin{enumerate}
    \item Se selecciona una vista representativa con el criterio de la
          sección~\ref{sec:seleccion_vista}.
    \item Se envía esa única imagen a Molmo2 con un prompt que solicita
          describir el objeto, y señalar puntos sobre el eje o plano de
          simetría en esa misma imagen.
    \item En base a la descripción del objeto y la indicación de los puntos
          de interés, se realiza un envío multivista en donde se pide hacer
          pointing a cada imagen.
\end{enumerate}

A diferencia de una lectura inicial de este flujo como línea base de cobertura
mínima, el paso 3 lo integra al mismo marco multivista que los Flujos A y B
($N \in \{1, 6, 14, 26\}$): la descripción y los puntos semilla obtenidos en el
paso 2 se usan como contexto para un prompt multivista aumentado, análogo al
del Flujo B pero informado además por ubicaciones concretas de puntos —no solo
por texto descriptivo.

\subsubsection{Selección de la vista de descripción}
\label{sec:seleccion_vista}

La selección aleatoria pura de entre las 114 vistas disponibles presenta un riesgo
concreto: vistas con elevación extrema ($|\text{el}| > 70°$) son
geométricamente degeneradas para la mayoría de objetos. Una vista cenital
($\text{el} \approx +89°$) proyecta el eje vertical de un objeto axial como un
punto y el plano de un objeto planar como una línea, produciendo descripciones
pobres que contaminarían el prompt de pointing.

Con 114 vistas distribuidas uniformemente en la esfera de Fibonacci,
aproximadamente el 10--15\% caen en elevaciones extremas ($|\text{el}| > 70°$),
lo que supone un riesgo no despreciable con muestreo aleatorio puro.

\paragraph{Estrategia adoptada: aleatoriedad filtrada por elevación.}
Se restringe el muestreo a las vistas con elevación en el rango
$\text{el} \in (-60°, +60°)$, que en la distribución Fibonacci corresponden a
aproximadamente 90--95 candidatas. Dentro de ese subconjunto se muestrea de forma
aleatoria con semilla fija derivada del identificador del objeto:

\begin{equation}
    s_{\text{obj}} = \mathrm{MD5}(\texttt{object\_id}) \bmod 2^{31}
    \label{eq:seed}
\end{equation}

El uso de una semilla por objeto (en lugar de una semilla global) garantiza que:
(i) la vista de descripción sea reproducible para cada objeto específico, (ii)
distintos objetos usen vistas distintas —evitando que todos sean descritos desde
la misma perspectiva—, y (iii) el experimento sea exactamente replicable.

\newpage
\subsection{Proyección 3D por \textit{ray casting}}
\label{sec:raycasting}

Las coordenadas 2D devueltas por Molmo2 se proyectan al espacio 3D mediante
\textit{ray casting} sobre la malla del objeto.

\paragraph{Paso 1: Conversión de coordenadas Molmo2 a NDC.}
Las coordenadas $(x, y) \in [0, 1000]^2$ se convierten a coordenadas normalizadas
de dispositivo (NDC, rango $[-1,1]$):

\begin{align}
    \text{ndc}_x &= \frac{x}{1000}\cdot 2 - 1 \label{eq:ndc_x} \\[4pt]
    \text{ndc}_y &= 1 - \frac{y}{1000}\cdot 2 \label{eq:ndc_y}
\end{align}

donde $\text{ndc}_x = -1$ es el borde izquierdo, $+1$ el derecho,
$\text{ndc}_y = +1$ el borde superior y $-1$ el inferior. Esta conversión es
independiente de la resolución efectiva de la imagen.

\paragraph{Paso 2: Construcción del rayo en espacio mundial.}
Bajo la convención de PyTorch3D ($p_{\text{cam}} = p_{\text{world}}\,R + T$),
el centro de la cámara en coordenadas mundiales es $\mathbf{O} = -R\,T$.
La dirección del rayo en espacio de cámara es:

\begin{equation}
    \mathbf{d}_{\text{cam}} =
    \begin{pmatrix}
        \text{ndc}_x \cdot \tan\!\left(\tfrac{\text{FoV}}{2}\right) \\[3pt]
        \text{ndc}_y \cdot \tan\!\left(\tfrac{\text{FoV}}{2}\right) \\[3pt]
        1
    \end{pmatrix}
    \label{eq:dir_cam}
\end{equation}

transformada a coordenadas mundiales y normalizada:

\begin{equation}
    \mathbf{d} = \frac{R\,\mathbf{d}_{\text{cam}}}{\|R\,\mathbf{d}_{\text{cam}}\|}
    \label{eq:dir_world}
\end{equation}

\paragraph{Paso 3: Intersección con la malla.}
El rayo $(\mathbf{O}, \mathbf{d})$ se intersecta contra la malla mediante Trimesh.
Se retiene la intersección más cercana; si no existe, el punto se descarta
(\textit{miss}).

\subsubsection{Variante: retroproyección por parche}
\label{sec:patchified}

La retroproyección píxel a píxel exacta es sensible a ángulos rasantes: cuando el
rayo llega casi tangente a la superficie, pequeños errores de localización en
píxeles producen grandes saltos en el punto 3D resultante. Como alternativa, se
propone evaluar una retroproyección por parche de $h \times h$ píxeles centrado
en el punto predicho:

\begin{equation}
    \mathbf{p}_{3D}^{\text{patch}} =
    \frac{1}{|S_{h}|}
    \sum_{(p,q)\,\in\, S_{h}}
    \mathrm{Intersect}\!\bigl(r_{p,q},\, \mathcal{M}\bigr)
    \label{eq:patch}
\end{equation}

donde $S_{h}$ es el conjunto de píxeles del parche $h\times h$ cuyos rayos
intersectan la malla $\mathcal{M}$, y el resultado es el promedio de todos esos
puntos 3D. El promediado estabiliza el resultado ante variaciones locales de
curvatura y errores de localización del VLM.

Se evaluará $h \in \{3, 5\}$ como tamaños de parche, comparando contra la
retroproyección exacta ($h=1$).

\subsection{Consolidación de la nube de puntos}
\label{sec:clustering}

Múltiples vistas pueden generar puntos 3D que corresponden a la misma región
física del objeto. Para reducir esa redundancia espacial se evalúan dos
estrategias de fusión:

\subsubsection{Clustering greedy por centroide (método actual)}

Los puntos se procesan en orden de llegada. Cada punto nuevo se asigna al primer
cluster cuyo centroide se encuentre dentro de un umbral $\tau_c = 0{,}05 \cdot
\delta_{\text{bbox}}$ (diagonal de la caja delimitadora de la nube de puntos);
si no existe tal cluster, se crea uno nuevo. El centroide se actualiza como media
incremental:

\begin{equation}
    \mathbf{c}_k \leftarrow \mathbf{c}_k +
    \frac{\mathbf{p} - \mathbf{c}_k}{n_k + 1}
\end{equation}

El método es $\mathcal{O}(n \cdot K)$ en número de clusters $K$ y asigna
\textit{siempre} cada punto a algún cluster —no existe concepto de outlier.

\subsubsection{HDBSCAN}

Como alternativa basada en densidad, se propone evaluar HDBSCAN
(\textit{Hierarchical Density-Based Spatial Clustering of Applications with
Noise})~\cite{campello2013hdbscan}. A diferencia del método greedy, HDBSCAN:

\begin{itemize}
    \item Detecta clusters como regiones de alta densidad sin requerir el número
          de clusters como parámetro.
    \item Clasifica explícitamente como \textit{ruido} los puntos que no
          pertenecen a ningún cluster denso, eliminándolos antes del ajuste SVD.
    \item Su parámetro principal es \texttt{min\_samples}: el número mínimo de
          puntos para que una región sea considerada densa.
\end{itemize}

La diferencia más relevante para este pipeline es el \textbf{manejo de outliers}:
mientras el clustering greedy siempre asigna cada punto a algún cluster (y delega
el rechazo de outliers al RANSAC posterior), HDBSCAN puede descartar
explícitamente puntos aislados —como los producidos por alucinaciones de Molmo2
que superan el \textit{hit} del ray casting pero caen en posiciones alejadas del
resto de la nube. Ambos métodos son estrategias distintas aplicadas en etapas
distintas: la combinación RANSAC + clustering greedy no es necesariamente inferior
a HDBSCAN solo.

Se evaluará \texttt{min\_samples} $\in \{2, 3, 5\}$ comparando contra el
clustering greedy y contra el caso sin clustering.

\newpage
\subsection{Detección y validación de simetrías}

Sobre la nube de puntos 3D consolidada se aplica un procedimiento de tres etapas
para estimar los parámetros de la simetría dominante del objeto.

\paragraph{Etapa 1 — Filtro de consenso (RANSAC).}
Se aplica RANSAC para descartar puntos atípicos (\textit{outliers}) provenientes
de detecciones erróneas del VLM. El umbral de inlier se fija en:

\begin{equation}
    \tau = 0{,}05 \cdot \delta_{\text{bbox}}
\end{equation}

donde $\delta_{\text{bbox}}$ es la diagonal de la caja delimitadora de la nube
de puntos 3D.

\begin{itemize}
    \item \textbf{Simetría axial}: cada iteración muestrea 2 puntos y define
          una recta candidata. La distancia de $\mathbf{q}$ a la recta es:
          \begin{equation}
              \delta = \left\|\,(\mathbf{q} - \mathbf{p}_i)
                          - \bigl[(\mathbf{q}-\mathbf{p}_i)\cdot\hat{\mathbf{d}}\bigr]
                          \hat{\mathbf{d}}\,\right\|
          \end{equation}

    \item \textbf{Simetría planar}: cada iteración muestrea 3 puntos y define un
          plano candidato con normal
          $\hat{\mathbf{n}} = (\mathbf{p}_2-\mathbf{p}_1)\times(\mathbf{p}_3-\mathbf{p}_1)$.
          La distancia de $\mathbf{q}$ al plano es:
          \begin{equation}
              \delta = \bigl|(\mathbf{q} - \mathbf{p}_1)\cdot\hat{\mathbf{n}}\bigr|
          \end{equation}
\end{itemize}

En ambos casos se clasifican como inliers los puntos con $\delta < \tau$. Tras
\textbf{1.000 iteraciones}, se selecciona el modelo candidato con mayor número de
inliers; si el conjunto ganador contiene menos de 2 inliers, se devuelven todos
los puntos como fallback.

\paragraph{Etapa 2 — Ajuste SVD.}
Sobre el conjunto de puntos seleccionado, centrado en su centroide $\bar{\mathbf{p}}$,
se aplica SVD: $P - \bar{\mathbf{p}} = U\,\Sigma\,V^\top$.

\begin{itemize}
    \item \textbf{Simetría axial}: $\hat{\mathbf{d}} = V^\top[0,:]$
          (primera componente principal, máxima varianza).
    \item \textbf{Simetría planar}: $\hat{\mathbf{n}} = V^\top[-1,:]$
          (última componente principal, mínima varianza).
\end{itemize}

El centroide $\bar{\mathbf{p}}$ se adopta como punto de origen de la simetría.

\paragraph{Etapa 3 — Verificación geométrica (SDE).}
Se muestrea una submuestra de hasta $N_{\text{SDE}} = 1.000$ vértices de la malla
y se aplica la transformación de simetría estimada.

Para \textbf{simetría axial} (reflexión de $180°$ respecto al eje):
\begin{equation}
    \mathbf{v}' = 2\Bigl(\mathbf{o} +
        \bigl[(\mathbf{v} - \mathbf{o})\cdot\hat{\mathbf{d}}\bigr]
        \hat{\mathbf{d}}\Bigr) - \mathbf{v}
    \label{eq:sde_axis}
\end{equation}

Para \textbf{simetría planar} (reflexión respecto al plano):
\begin{equation}
    \mathbf{v}' = \mathbf{v} -
        2\bigl[(\mathbf{v} - \mathbf{o})\cdot\hat{\mathbf{n}}\bigr]\hat{\mathbf{n}}
    \label{eq:sde_plane}
\end{equation}

Para cada vértice transformado $\mathbf{v}'$ se busca su vecino más cercano en la
submuestra mediante un árbol KD:
\begin{equation}
    \text{SDE} = \frac{1}{|V|}\sum_{\mathbf{v} \in V}
        \frac{\|\mathbf{v}' - \mathrm{NN}(\mathbf{v}')\|}{\delta_{\text{mesh}}}
    \label{eq:sde}
\end{equation}

\textbf{Criterio de aceptación}: $\text{SDE} > 0{,}05 \Rightarrow$
\texttt{accepted = false} (umbral fijado empíricamente).

La combinación de estas etapas produce cuatro estimaciones por objeto y número
de vistas:

\begin{center}
\begin{tabular}{lcc}
\toprule
 & Sin SDE & Con SDE \\
\midrule
Sin RANSAC & \texttt{svd} & \texttt{svd\_sde} \\
Con RANSAC & \texttt{ransac\_svd} & \texttt{ransac\_svd\_sde} \\
\bottomrule
\end{tabular}
\end{center}

\newpage
\subsection{Evaluación cuantitativa}

El desempeño del sistema se evalúa comparando los parámetros predichos contra las
anotaciones de referencia:

\paragraph{Error angular.}
\begin{equation}
    \varepsilon_\theta = \arccos\!\bigl(\,|\hat{\mathbf{d}}_{\text{pred}}
        \cdot \hat{\mathbf{d}}_{\text{gt}}|\,\bigr) \in [0°,\,90°]
\end{equation}

Para simetría planar: \textit{best-match} contra todos los planos de referencia
disponibles.

\paragraph{Error de traslación.}
Distancia punto-recta (axial) o punto-plano (planar) del origen predicho al GT,
normalizadas por $\delta_{\text{mesh}}$.

\paragraph{Precisión bajo umbral angular.}
Fracción de objetos con $\varepsilon_\theta < \{5°, 10°, 15°\}$.

\paragraph{AUC angular.}
\begin{equation}
    \text{AUC}_\theta =
    \frac{1}{45}\int_0^{45}\text{Precision}(t)\,dt
\end{equation}

Los objetos sin predicción reciben imputación $\varepsilon_\theta = 90°$.

\paragraph{SDE de evaluación — exclusivo simetría planar.}
Calculado sobre la totalidad de los vértices de la malla:
\begin{equation}
    \text{SDE}_{\text{eval}} =
    \frac{1}{|V|}\sum_{\mathbf{v}\in V}
    \frac{2\,\bigl|(\mathbf{v}-\mathbf{o})\cdot\hat{\mathbf{n}}\bigr|}
         {\delta_{\text{mesh}}}
    \label{eq:sde_eval}
\end{equation}

\paragraph{Precisión y AUC bajo SDE.}
\begin{equation}
    \text{AUC}_{\text{SDE}} =
    \frac{1}{0{,}10}\int_0^{0{,}10}
        \text{Precision}_{\text{SDE}}(t)\,dt
\end{equation}

con umbrales de $\{0{,}01,\, 0{,}02\}$.

Las métricas angulares se calculan sobre los 100 objetos evaluados (imputando
$90°$ a los sin predicción); el error de traslación y el SDE solo sobre los
$n_{\text{obj}}$ objetos con predicción aceptada.

\newpage
\subsection{Diseño experimental}
\label{Experimentos}

Los experimentos se organizan en tres flujos de procesamiento independientes,
dos variantes de clustering y dos variantes de retroproyección. Las comparaciones
se realizan siempre bajo las mismas métricas (§3.6) y sobre el mismo subconjunto
de 100 objetos, a fin de mantener la interpretabilidad de los resultados.

\subsubsection{Flujo A — Pointing directo (línea base)}

El modelo recibe las $N$ vistas y un prompt de pointing sin contexto adicional.
Este flujo replica y extiende los experimentos ya realizados:

\begin{itemize}
    \item \textbf{Prompts evaluados}: 6 diseños originales (v00--v05) $\times$
          2 tipos de simetría, ya completados.
    \item \textbf{Prompts mejorados}: 6 versiones v\_1 (v00\_1--v05\_1) $\times$
          2 tipos de simetría, ya completados.
    \item \textbf{Condiciones}: resolución 224~px, iluminación \textit{flat},
          $N \in \{1, 6, 14, 26\}$ vistas, 4 métodos de ajuste.
\end{itemize}

\subsubsection{Flujo B — Pointing con descripción (descripción prepend)}

El modelo recibe la misma configuración multivista que en el Flujo A, pero el
prompt se aumenta con una descripción semántica del objeto generada previamente
desde una vista de descripción seleccionada según el criterio de la
sección~\ref{sec:seleccion_vista}:

\begin{itemize}
    \item \textbf{Vista de descripción}: aleatoria con semilla por objeto
          (ecuación~\eqref{eq:seed}), filtrada a elevación
          $\text{el} \in (-60°, +60°)$.
    \item \textbf{Prompt}: descripción libre de Molmo2 + prompt de pointing base
          (best performer del Flujo A).
    \item \textbf{Condiciones}: resolución 224~px, iluminación \textit{flat},
          $N \in \{1, 6, 14, 26\}$ vistas, 4 métodos de ajuste.
\end{itemize}

Este flujo permite evaluar si el contexto semántico mejora la localización de
los puntos de simetría respecto al Flujo A.

\subsubsection{Flujo C — Descripción y pointing integrados}

El modelo recibe inicialmente una sola vista para la descripción del objeto y
de los puntos a encontrar. Posteriormente se le pasan $N$ vistas y un prompt
general para simetría axial y planar que solicita el pointing para los puntos
señalados:

\begin{itemize}
    \item \textbf{Vista}: misma selección que en Flujo B (aleatoria filtrada por
          elevación, semilla por objeto).
    \item \textbf{Prompt}: un único mensaje inicial que solicita describir el
          objeto y señalar los puntos sobre el eje/plano de simetría en esa
          imagen. Posteriormente se incluye dicha información en otro prompt
          general que tiene como objetivo realizar pointing.
    \item \textbf{Resultado}: múltiples puntos 3D por objeto.
    \item \textbf{Condiciones}: resolución 224~px, iluminación \textit{flat},
          $N \in \{1, 6, 14, 26\}$ vistas, 4 métodos de ajuste.
\end{itemize}

Este flujo cuantifica cuánto aporta, sobre la base multivista común a los tres
flujos, el contexto adicional de una descripción y puntos semilla obtenidos de
una única vista de referencia.

\subsubsection{Variantes de clustering}

Para cada flujo se evaluarán dos estrategias de consolidación de la nube de
puntos (§\ref{sec:clustering}):

\begin{center}
\begin{tabular}{lll}
\toprule
Variante & Método & Parámetros \\
\midrule
C1 & Sin clustering & — \\
C2 & Greedy por centroide & $\tau_c = 0{,}05 \cdot \delta_{\text{bbox}}$ \\
C3 & HDBSCAN & \texttt{min\_samples} $\in \{2, 3, 5\}$ \\
\bottomrule
\end{tabular}
\end{center}

\subsubsection{Variantes de retroproyección}

Para el mejor prompt de los Flujos A y B se evaluará la retroproyección por
parche (§\ref{sec:patchified}) como alternativa a la retroproyección exacta:

\begin{center}
\begin{tabular}{lll}
\toprule
Variante & Método & Parámetros \\
\midrule
R1 & Ray casting exacto & $h = 1$ (píxel único) \\
R2 & Retroproyección por parche & $h \in \{3, 5\}$ \\
\bottomrule
\end{tabular}
\end{center}

\subsubsection{Resumen del plan experimental}

\begin{center}
\begin{tabular}{llccc}
\toprule
Flujo & Descripción & Prompts & $N$ vistas & Clustering \\
\midrule
A & Pointing directo — originales      & v00--v05   & 1,6,14,26 & C1, C2 \\
A & Pointing directo — mejorados       & v00\_1--v05\_1 & 1,6,14,26 & C1, C2 \\
B & Pointing + descripción             & best(A)    & 1,6,14,26 & C1, C2, C3 \\
C & Descripción y pointing integrados   & nuevo      & 1,6,14,26 & C1, C2, C3 \\
\midrule
  & Retroproyección por parche (best)  & best(A/B)  & best($N$) & best(C) \\
\bottomrule
\end{tabular}
\end{center}

Los experimentos de Flujo A con prompts originales y mejorados están completados
($\checkmark$). Los Flujos B y C, las variantes HDBSCAN (C3) y la retroproyección
por parche (R2) constituyen el trabajo experimental pendiente.

\subsection{Comparación con métodos tradicionales basados en geometría}

Los resultados obtenidos se contrastarán con el método propuesto por Aguirre y
Sipiran~\cite{aguirre_sipiran}, que aborda la detección de simetría 3D mediante
retroproyección de características visuales con acceso directo a la geometría del
objeto. La comparación se realiza bajo las mismas métricas (§3.6) y sobre el
mismo subconjunto de objetos de ShapeNet, permitiendo cuantificar el costo, en
términos de precisión geométrica, de prescindir del acceso a la malla 3D durante
la extracción de puntos —sustituyéndola por un VLM operando exclusivamente sobre
observaciones 2D.

\newpage
\subsection*{Entorno de implementación}

El sistema fue implementado en Python~3.10 con las siguientes dependencias
principales: PyTorch~2.6 (CUDA~12.4), PyTorch3D~0.7.9 para el renderizado y la
cámara perspectiva, Trimesh para \textit{ray casting}, Transformers~4.57 +
Molmo2-8B para la inferencia del VLM, y NumPy/SciPy para álgebra lineal y
construcción del árbol KD. Los experimentos de inferencia se ejecutaron en GPU;
las etapas de \textit{ray casting}, ajuste SVD y evaluación son tareas de CPU
ejecutadas en paralelo mediante segmentación por objeto.

La correspondencia entre el sistema de coordenadas del renderizado y el del
archivo \texttt{.obj} original está garantizada por diseño: los parámetros $R$,
$T$ almacenados en \texttt{metadata\_all.json} son los mismos que se usan tanto
para el renderizado como para la reconstrucción del rayo en
\texttt{map\_to\_3d.py}, y la posición de cámara $\mathbf{C} = -RT$ fue validada
numéricamente verificando su coincidencia con el parámetro \textit{eye} de
\texttt{look\_at\_view\_transform}.
